
\section{An equational presentation of stochastic matrices}\label{sec:syntactic}    

\begin{table}[t]
\resizebox{\textwidth}{!}{%
$
\begin{array}{c}
\toprule
\begin{array}{c}
    \begin{array}{ccccc}
        {\diagp{U}\colon U \to U\oplus U} &
        {\bang{U} \colon U \to \zero} &
        {\sigma_{U, V}^{\piu} \colon U \piu V \to V \piu U} &
        {\codiag{U}\colon U \piu U \to U} &
        {\cobang{U} \colon \zero \to U}   
    \end{array}    
\\[0.3em]
    \begin{array}{ccccc}
        {id_\zero \colon \zero \to \zero} &
        {id_U \colon U \to U} &
        \inferrule{c \colon U \to V}{\tapeFunct{c}\colon U \to V} &
        \inferrule{\t \colon P \to Q \and \s \colon Q \to R}{\t ; \s \colon P \to R} &
        \inferrule{\t \colon P_1 \to Q_1 \and \s \colon P_2 \to Q_2}{\t \piu \s \colon P_1 \piu P_2 \to Q_1 \piu Q_2}       
    \end{array}
\end{array}
\\[0.5em]
\begin{array}{cc}
\midrule
\begin{array}{c}
(f;g);h=f;(g;h) \qquad id_P;f=f=f;id_Q\\
(f_1\piu f_2) ; (g_1 \piu g_2) = (f_1;g_1) \piu (f_2;g_2)
\end{array} 
&
\begin{array}{c}
id_{\zero}\piu f = f = f \piu id_{\zero} \qquad (f \piu g)\, \piu h = f \piu \,(g \piu h) \\
\sigma_{P, Q}; \sigma_{Q, P}= id_{P \piu Q} \qquad (\gen \piu id_R) ; \sigma_{Q, R} = \sigma_{P,R} ; (id_R \piu \gen)
\end{array}
\end{array}\\
\bottomrule
\end{array}
$
}
\caption{Typing rules (top) and axioms (bottom) for freely generated strict symmetric monoidal categories.}
\label{fig:freestricmmoncatax}
\end{table}

%\begin{table}
%\begin{center}
%\begin{tabular}{c}
%$(\id{ A}\piu \codiag{ A}) ; \codiag{ A} = (\codiag{ A}\piu \id{ A}) ; \codiag{ A} \qquad (\cobang{ A}\piu \id{ A}) ; \codiag{ A}  = \id{ A} \qquad \sigma_{ A, A};\codiag{ A}=\codiag{ A}$ \\

%$\cobang{\zero} = \id{\zero} \qquad\codiag{\zero} = \id{\zero} \qquad \cobang{A \piu W} = \cobang{A} \piu \cobang{W} \qquad \codiag{A \piu W} = (\id{A} \piu \sigma_{A,W} \piu \id{W}) ; (\codiag{A}\piu \codiag{W})$\\
%$\cobang{X};f =\cobang{Y} \qquad \codiag{X};f \mylabel{ax:codiagnat}{$\codiag{}$-nat}(f\piu f); \codiag{Y}$


%\end{tabular}
%\end{center}
%\caption{Additional axioms for freely generated strict fc categories.}\label{fig:freestrictfccat}
%\end{table}







% \begin{table}[ht!]
%     \begin{center}
%         \begin{tabular}{c|c}
%             \toprule
%             \begin{tabular}{l}
%                 % \toprule
%                 $\dl{P}{Q}{R} \colon P \per (Q\piu R)  \to (P \per Q) \piu (P\per R) \vphantom{\symmt{P}{Q}}$ \\
%                 \midrule
%                 $\dl{\zero}{Q}{R} \defeq \id{\zero} \vphantom{\symmt{P}{\zero} \defeq \id{\zero}}$ \\
%                 $\dl{U \piu P'}{Q}{R} \defeq (\id{U\per (Q \piu R)} \piu \dl{P'}{Q}{R});(\id{U\per Q} \piu \symmp{U\per R}{P'\per Q} \piu \id{P'\per R}) \vphantom{\symmt{P}{V \piu Q'} \defeq \dl{P}{V}{Q'} ; (\Piu[i]{\tapesymm{U_i}{V}} \piu \symmt{P}{Q'})}$ \\
%                 % \bottomrule
%             \end{tabular}
%             &
%             \begin{tabular}{l}
%                 % \toprule
%                 $\symmt{P}{Q} \colon P\per Q \to Q \per P$, with $P = \Piu[i]{U_i}$ \\
%                 \midrule
%                 $\symmt{P}{\zero} \defeq \id{\zero}$ \\
%                 $\symmt{P}{V \piu Q'} \defeq \dl{P}{V}{Q'} ; (\Piu[i]{\tapesymm{U_i}{V}} \piu \symmt{P}{Q'})$ \\
%                 % \bottomrule
%             \end{tabular}
%             \\[20pt]
%             \midrule
%             \multicolumn{2}{c}{
%                 \begin{tabular}{c|c|c}
%                     \begin{tabular}{l}
%                         % \toprule
%                         $\codiag{P} \colon P \per P \to P$
%                         \\
%                         \midrule
%                         $\codiag{\zero} \defeq \id{\zero} $
%                         \\
%                         $\codiag{U \piu P'} \defeq (\id{U} \piu \symmp{U}{P'} \piu \id{P'}) ; (\codiag{U} \piu \codiag{P'})$
%                         % \\
%                         % \bottomrule
%                     \end{tabular}
%                     &
%                     \begin{tabular}{l}
%                         % \toprule
%                         $\cobang{P} \colon \zero \to P
%                         \vphantom{\copier{U} \colon U \to U \per U}$
%                         \\
%                         \midrule
%                         $\cobang{\zero} \defeq \id{\zero} $
%                         \\
%                         $\cobang{U \piu P'} \defeq \cobang{U} \piu \cobang{P'}
%                         \vphantom{\cocopier{U \piu P'} \defeq \Tcocopier{U} \piu \bang{UP'} \piu (\dl{P'}{U}{P'} ; (\bang{P'U} \piu \cocopier{P'}))}$
%                         % \\
%                         % \bottomrule
%                     \end{tabular}
%                     &
%                     \begin{tabular}{l}
%                         % \toprule
%                         $f_P \colon P \to \piu^n P$ \\
%                         \midrule
%                         $f_\zero \defeq \id{\zero}$ \\
%                         $f_{U \piu P'} \defeq (f_U \piu f_{P'}) ; \Idl{\piu^n \uno}{U}{P'}$ 
%                         % \bottomrule
%                     \end{tabular}
%                 \end{tabular}
%             }
%             \\
%             \bottomrule
%         \end{tabular}
%         \caption{Syntactic sugar for $\algt$-tapes.}
%         \label{table:def syn sugar}
%     % }
%     \end{center}
% \end{table}


% \begin{equation}\label{eq:adjunction} \begin{tikzcd}
%     \SMC \ar[r,"U_1"] & \CAT \ar[r,bend left,"F_2"] \ar[r,draw=none,"\bot"description] & \Cat{FCC} \ar[l,bend left,"U_2"]
% \end{tikzcd}
% \end{equation}
% The functors $U_1$ and $U_2$ are the obvious forgetful functors. The functor $F_2$ is the left adjoint to $U_2$, and can be described as follows.

% \begin{definition}\label{def:strict fb freely generated by C}
% Let $\Cat{C}$ be a category. The  strict fc category freely generated by $\Cat{C}$, hereafter denoted by $F_2(\Cat{C})$, has as objects words of objects of $\Cat{C}$. Arrows are terms inductively generated by the following grammar, where $A,B$ and $c$ range over arbitrary objects and arrows of $\Cat{C}$,\begin{equation} 
%     \begin{array}{rcl}
%         f & ::=& \; \id{A} \; \mid \; \id{I} \; \mid \; \tapeFunct{c} \; \mid \; \sigma_{A,B}^{\perG} \; \mid \;   f ; f   \; \mid \;  f \perG f  \;  \mid \; \cobang{A}\; \mid \; \codiag{A} \\ 
%         \end{array}
% \end{equation}  
% modulo the axioms in  Table~\label{tab:tape typing and axioms}. Notice in particular the last two from the bottom right of Table\label{tab:tape typing and axioms}:
% \begin{equation}\label{ax:tape}\tag{Tape}
% \tape{\id{A}} = \id{A} \qquad \tape{c;d}= \tape{c} ; \tape{d}
% \end{equation}
% \end{definition}
% The assignment $\Cat{C} \mapsto F_2({\Cat{C}})$ easily extends to functors $H\colon \Cat{C} \to \Cat{D}$.
% The unit of the adjunction $\eta \colon Id_{\CAT} \Rightarrow F_2U_2$ is defined for each category $\Cat{C}$ as the functor ${\tapeFunct{\cdot} \; \colon \Cat{C} \to U_2F_2(\Cat{C})}$ which is the identity on objects and maps every arrow $c$ in $\Cat{C}$ into the arrow $\tapeFunct{c}$ of $U_2F_2(\Cat{C})$. Observe that $\tapeFunct{\cdot}$ is indeed a functor, namely an arrow in $\CAT$, thanks to the axioms \eqref{ax:tape}. We will refer hereafter to this functor as the \emph{taping functor}.


% \begin{lemma}\label{lemma:F2 left adjoint to U2}
% $F_2 \colon \CAT \to \FBC$ is left adjoint to $U_2 \colon \FBC \to \CAT$.
% \end{lemma}

% \smallskip

% AGGIUSTARE DA QUA IN POI
% \smallskip

So far, we have seen that given a category, one can first freely enrich it over $\Cat{PCA}$, and then obtain a convex biproduct category by taking its category of stochastic matrices.
Now, we give a syntactic description, in terms of generators and equations, of such a category.


For a category $\Cat{C}$, we consider terms generated by the following grammar
  
\begin{equation}\label{tapesGrammar}
\setlength{\arraycolsep}{3pt}
\begin{array}{rcccccccccccccccccccc}
\t & ::= & \diagp{U} & \mid & \bangp{U} & \mid & \tapeFunct{c} & \mid & \cobang{U} & \mid & 
\codiag{U} & \mid & \id{U} & \mid & \id{\zero} & \mid & \sigma_{U,V}^{\piu} & \mid & \t ; \t & \mid & \t \piu \t
\end{array}
\end{equation}
%
%
%\begin{equation}\label{tapesGrammar}
%    \begin{tabular}{rc ccccccccccccccccccccc}\setlength{\tabcolsep}{0.0pt}
%        $\t$ & ::= & $\id{A}$ & $\!\!\! \mid \!\!\!$ & $ \id{\zero} $ & $\!\!\! \mid \!\!\!$ & $\Br{+_p}_A$  & $\!\!\! \mid \!\!\!$  & $\Br{\star}_A$  & $\!\!\! \mid \!\!\!$  & $ \tapeFunct{c} $ & $\!\!\! \mid \!\!\!$ & $ \sigma_{A,B}^{\piu} $ & $\!\!\! \mid \!\!\!$ & $   \t ; \t   $ & $\!\!\! \mid \!\!\!$ & $  \t \piu \t  $ & $\!\!\! \mid \!\!\!$  & $\cobang{A}$ & $\!\!\! \mid \!\!\!$ & $\codiag{A}$ & $\!\!\! \mid \!\!\!$ 
%    \end{tabular}
%\end{equation}  
where $p\in (0,1)$, $U,V \in \ob{\Cat{C}}$, and $c$ is an arrow in $\Cat{C}$. Terms are typed according to the rules at the top of Table~\ref{fig:freestricmmoncatax}: each type is an arrow $P \to Q$ where $P,Q\in \ob{\Cat{C}}^*$ are regarded as sums of objects of $\Cat{C}$. As expected, we consider only those terms that are typable. For arbitrary $P \in \ob{\Cat{C}}^*$, we define
$\diagp{P},\bangp{P},\cobang{P},\codiag{P}$ as in \eqref{eq:indmoncopca}. Analogous inductive definitions give us $\id{P}$ and $\symmp{P}{Q}$. %\marginpar{Think later about the types} %Terms are taken modulo the axioms in BLA BLA

The category $\CatTapeC$ has as set of objects $\ob{\Cat{C}}^*$. Arrows in $\CatTapeC[P,Q]$ are terms modulo the axioms of natural and coherent monoids, co-pcas, strict symmetric monoidal categories (respectively in Tables~\ref{fig:freestrictfccat}, \ref{fig:freecopcacat} and~\ref{fig:freestricmmoncatax} bottom) and the following two axioms.
\begin{equation}\label{eq:TapeFunctAxioms}
\tapeFunct{\id{P}}=\id{P} \qquad \tapeFunct{c;d} = \tapeFunct{c}; \tapeFunct{d} \tag{Tape}
\end{equation}
Identities, composition, symmetries and monoidal product are defined as for terms. It is thus immediate to see that $(\CatTapeC,\oplus, \zero)$ forms a symmetric monoidal category that, like in Definition~\ref{definition:structural}, is equipped with natural and coherent monoid and co-pca structures. To prove that $\CatTapeC$ is a convex biproduct category, by means of Proposition~\ref{prop:A}, we need to prove that projections are jointly monic.



%. Moreover, Lemma~\ref{lemma:enrichmentpca}, implies that such category is enriched over $\Cat{PCA}$ and  for all objects $P,Q$ and arrows $f,g\colon P \to Q$
%\begin{equation}\label{eq:enrichmentTC}
%f +_p g \defeq\   \diagp{P} ; (f\oplus g) ; \codiag{Q} \qquad \star_{P,Q} \defeq \, \bangp{P} ; \cobang{Q}
%\end{equation}
%By naturality of $\cobang{}$ and $\bang{}$, $\zero$ is both initial and final object. By the Fox theorem \cite{fox1976coalgebras}, $\oplus$ is a coproduct with injections $\id{P_1}\oplus \cobang{P_2} \colon P_1 \to P_1\oplus P_2$ and $\cobang{P_1}\oplus \id{P_2} \colon P_2 \to P_1\oplus P_2$. %Most importantly, $\oplus$ is a convex product with projections $\id{P_1}\oplus \bang{P_2} \colon P_1 \oplus P_2 \to P_1$ and $\bang{P_1}\oplus \id{P_2} \colon P_1\oplus P_2 \to  P_2$.

%
%
%
%
%Our main technical effort consists in proving that $\oplus$ is a convex product with projections $\id{P_1}\oplus\, \bang{P_2} \colon P_1 \oplus P_2 \to P_1$ and $\bang{P_1}\oplus \id{P_2} \colon P_1\oplus P_2 \to  P_2$. 

We crucially rely on the fact that every object $P$ of $\CatTapeC$ is of the form $\bigoplus_{i=1}^nU_i$ where $\oplus$ is a coproduct by Lemma~\ref{lemma:coproducts}. Thanks to the universal property of coproducts, any $\t\colon\bigoplus_{j=1}^mV_j \to \bigoplus_{i=1}^nU_i$ is the copairing $[\t_1, \dots, \t_m]$ for $\t_j=\iota_j; \t$. One can thus restrict to considering the case of arrows of type $ V_j \to \bigoplus_{i=1}^nU_i$. These arrows enjoy a convenient normal form, illustrated in the next lemma. First, for all $n\in \mathbb{N}$ and $\vec{p}=p_1,\dots ,p_n$ with $p_i\in[0,1]$ such that $\sum_{i=1}^n p_i\leq 1$, we inductively define $\diagpn{\vec{p}\;}{U}{n} \colon U \to \bigoplus_{i=1}^nU$ as follows
\begin{equation}\label{eq: diagpn cap equational presentation}
\diagpn{\vec{p}\;}{U}{0}\defeq \bang{U} \qquad \diagpn{\vec{p}\;}{U}{n+1}\defeq\, \diagpn{\;p_{1}}{U}{}; (\id{U} \oplus \diagpn{\;\vec{q}\;}{U}{n})
\end{equation}
where $\vec{q}=q_1,\dots q_{n}$ for $q_i=0$ if $p_{1}= 1$ and $q_i=\frac{p_{i+1}}{1-p_{1}}$ otherwise. Note that the arrow in \eqref{eq: diagpn cap equational presentation} coincides with the one in \eqref{diagpq generalizzata} for $n=2$.
%Given $n$ arrows $t_i\colon U \to U_i$, simple computations confirm that $(\diagpn{\;\vec{p}\;}{U}{n}; \bigoplus_{i=1}^n \t_i) ; \pi_i = p_i\cdot \t_i$. Thus, $\diagpn{\;\vec{p}\;}{U}{n}; \bigoplus_{i=1}^n \t_i \colon U \to \bigoplus_{i=1}^nU_i$ is a mediating arrow satisfying the constraints of $n$-ary convex products. The hard part consists in proving its uniqueness. We rely on three key properties of $\CatTapeC$:
\begin{lemma}\label{decomposition}
For all $\t\colon U \to \bigoplus_{i=1}^n Q_i$, there exist $\vec{p}=p_1,\dots ,p_n$ and $t_i\colon U \to Q_i$ such that \[\t=\,\diagpn{\;\vec{p}\;}{U}{n}; \bigoplus_{i=1}^n \t_i\text{.}\]
\end{lemma}

The second key property is that the enrichment, defined as in \eqref{eq:enrichment}, is cancellative.

\begin{lemma}\label{cancellativity}[Cancellativity]
For all $r\in (0,1)$, for all $\s,\t \colon P \to Q$, 
if $r\cdot \s = r\cdot \t$ then $\s = \t$.
\end{lemma}

By means of the two lemmas above one can easily conclude that projections, defined as in \eqref{eq:defpi}, are jointly monic.

\begin{lemma}\label{lemma:TC jmono}
$\CatTapeC$ has jointly monic projections.
\end{lemma}
%
%
%

\begin{theorem}\label{thm:TCconvexbiproductcategory}
$\CatTapeC$ is a convex biproduct category. In particular, for all $\vec{p}=p_1,\dots ,p_n$ and $\t_i\colon U \to Q_i$, $\langle \t_1, \dots \t_n\rangle_{\vec{p}} = \diagpn{\;\;\vec{p}}{U}{n}; \bigoplus_{i=1}^n \t_i$.
\end{theorem}
\begin{proof}
    By construction, every object in $\CatTapeC$ is a natural and coherent comonoid and co-pca object, and Lemma~\ref{lemma:TC jmono} implies that the projections are jointly monic. Hence, $\CatTapeC$ is a convex biproduct category by Proposition~\ref{prop:A}. The last part of the statement follows from Lemma~\ref{decomposition} and uniqueness of mediating arrows for convex products.
\end{proof}
The assignment $\Cat{C} \mapsto \CatTapeC$ gives rise to a functor $\CatTapeFUN\colon \Cat{Cat} \to \Cat{CBCat}$ which is left adjoint to the forgetful functor  $U\colon \Cat{CBCat} \to \Cat{Cat} $.

\begin{theorem}\label{thm:syntacticadjunction}
$\CatTapeFUN\colon \Cat{Cat} \to \Cat{CBCat}$ is left adjoint to $U\colon \Cat{CBCat} \to \Cat{Cat} $.
\end{theorem}
\begin{proof}
For a functor $F\colon \Cat{C} \to \Cat{D}$, $\CatTapeF \colon \CatTapeC \to \CatTapeD$ is defined on objects as $\CatTapeF(\bigoplus_{i=1}^nU_i) = \bigoplus_{i=1}^nF(U_i)$.
%\[\CatTapeF(U) \defeq F(U) \qquad \CatTapeF(\zero)\defeq\zero \qquad \CatTapeF(P\oplus Q) \defeq   \CatTapeF(P)\oplus \CatTapeF(Q) \]
%where as usual $U$ is an object of $\Cat{C}$.
On arrows, it is defined inductively as follows.
\[
\setlength{\arraycolsep}{2pt} % riduce lo spazio tra colonne
\begin{array}{rcl rcl}
  \CatTapeF(\t_1 ; \t_2)      &\defeq& \CatTapeF(\t_1);\CatTapeF(\t_2) &
  \CatTapeF(\t_1 \piu \t_2)   &\defeq& \CatTapeF(\t_1)\oplus\CatTapeF(\t_2) \\
  
  \CatTapeF(\diagp{P})        &\defeq& \diagp{\CatTapeF(P)} &
  \CatTapeF(\codiag{P})       &\defeq& \codiag{\CatTapeF(P)} \\
  
  \CatTapeF(\,\bangp{P})        &\defeq& \bangp{\CatTapeF(P)} &
  \CatTapeF(\cobang{P})       &\defeq& \cobang{\CatTapeF(P)} \\
  
  \CatTapeF(\id{U})           &\defeq& \id{\CatTapeF(U)} &
  \CatTapeF(\id{\zero})       &\defeq& \id{\CatTapeF(\zero)} \\
  
  \CatTapeF(\sigma_{P,Q}^{\piu}) &\defeq& \sigma_{\CatTapeF(P),\CatTapeF(Q)}^{\piu} &
  \CatTapeF(\tapeFunct{c})    &\defeq& \tapeFunct{F(c)}
\end{array}
\]


Since by definition $\CatTapeF$ is a monoidal functor preserving co-pcas and monoids, by Proposition~\ref{prop: monoidal functors}, it is a morphism of convex biproduct categories.
%
Thus, we have a functor $\CatTapeFUN\colon \Cat{Cat} \to \Cat{CBCat}$. Below, we prove the adjunction.

For all categories $\Cat{C}$, the unit of the adjunction $\eta_{\Cat{C}} \colon \Cat{C} \to \CatTapeC$ is the identity-on-objects functor mapping each arrow $c\colon U\to V$ in $\Cat{C}$ into $\tapeFunct{c}$. The axioms in \eqref{eq:TapeFunctAxioms} force $\eta_{\Cat{C}}$ to be a functor. Naturality of the unit is straightforward.

Now, take a convex biproduct category $\Cat{D}$ and consider a functor $F\colon \Cat{C} \to U(\Cat{D})$. By Proposition~\ref{lemma: copca objects in convbicat}, $\Cat{D}$ is a symmetric monoidal category where every object $X$  carries a  co-pca $(\diagp{X},\bang{X})$ and a natural coherent monoid $(\codiag{X},\cobang{X})$. One can use these structures to define $F^\sharp\colon \CatTapeC \to \Cat{D}$ inductively in the same way as  $\CatTapeF$: e.g., $F^\sharp(\diagp{P})\defeq \,\diagp{F^\sharp(P)}$. The base case $F^\sharp(\tapeFunct{c})=F(c)$ ensures that $\eta_{\Cat{C}}; F^\sharp = F$.
Since by definition $F^\sharp$ is a strict symmetric monoidal functor preserving co-pcas and monoids, by Proposition~\ref{prop: monoidal functors}, it is a morphism of convex biproduct categories. 

For uniqueness, take a morphism of convex biproduct categories $H \colon \CatTapeC \to \Cat{D}$ such that $\eta_{\Cat{C}}; H = F$, i.e.,  $H(\tapeFunct{c})=F(c)$. By Proposition~\ref{prop: monoidal functors}, $H$  preserves co-pcas and monoids. Thus $H=F^\sharp$.
\end{proof}

\begin{corollary}\label{cor:isotapematrices}
For all categories $\Cat{C}$, $\CatTapeC$ is isomorphic to $\stmat{\Cat{C}^+}$ in $\Cat{CBCat}$.
\end{corollary}
\begin{proof}
By composition of adjoints, their uniqueness and Theorems~\ref{thm:freeenriched}, \ref{thm:matfree} and \ref{thm:syntacticadjunction}.
\end{proof}
The isomorphism $\CatTapeC \to \stmat{\Cat{C}^+}$ maps $\tapeFunct{c}$ into the 
$1\times 1$ matrix $\begin{pmatrix}
   1 \cdot c
\end{pmatrix}$ and  $\diagp{U}$, $\bang{U}$, $\cobang{U}$, $\codiag{U}$ into the matrices defined in \eqref{eq:matmonpca}.  Compositions and sums of arrows in $\CatTapeC$ %\marginpar{Please Cipriano Check}
are mapped into multiplications and direct sums of matrices as defined in \eqref{eq:matrixmult} and \eqref{eq:matsmc}. The same applies to identities and symmetries. For instance, $(\diagpX{\frac{1}{2}\;\,}{U} \oplus \id{U}) ; (\tapeFunct{a}\oplus \tapeFunct{ab} \oplus \tapeFunct{c}); (\id{U}\oplus \symmp{U}{U}); (\codiag{U}\oplus \id{U})$ is mapped into the matrix $M$ of Example~\ref{ex:matrices}. 

The characterisation of $\stmat{\Cat{C}^+}$ by means of generators and equations provided by the above corollary, paves the way to study further equational theories (such as the one in Section~\ref{sec:probbooltapes}). The following result guarantees that the category obtained by quotienting $\CatTapeC$ by additional axioms is still a convex biproduct category.
%
%We conclude with a result that will be useful in Section \ref{sec:probboolcircuits}.


\begin{proposition}\label{cor: quotient category is convex biproduct }
Let $\Cat{C}$ be a category and let $\sim$ be a congruence relation (w.r.t.\ $;$ and $\oplus$) on $\CatTapeC$ which is cancellative: for every $p\in (0,1)$, $p\cdot f\sim p\cdot g$ then $f\sim g$. Let $\CatTapeC_\sim$ be the category obtained as the quotient of $\CatTapeC$ by $\sim$ and $Q_\sim \colon \CatTapeC\to \CatTapeC_\sim$ be the functor mapping each arrow to its $\sim$-equivalence class.
Then $\CatTapeC_\sim$ is a convex biproduct category and $Q_\sim$ is a morphism of convex biproduct categories.
\end{proposition}